% ================================================================
% MODULE: Skills
% ================================================================
% AI INSTRUCTIONS:
% Fill placeholders with job-relevant skills from asset library.
%
% Placeholders to fill:
%   - {{SKILLS_SECTION_TITLE}} (optional: customize section title if needed)
%   - {{TECHNICAL_SKILLS_BULLETS}} (3-6 category bullets)
%
% SECTION TITLE CUSTOMIZATION:
%   - Default: "Technical \& Computational Skills"
%   - For Software/ML jobs: "Technical Skills" or "Programming \& Data Science"
%   - For Research jobs: "Technical \& Computational Skills" or "Research \& Computing Skills"
%   - For Industry jobs: "Technical \& Engineering Skills"
%   - Keep it concise and relevant to {{JOB_KEYWORDS}}
%
% CATEGORIZATION RULES:
%   - Group skills by category (e.g., Programming, Tools, Domain Knowledge)
%   - Each category = 1 bullet point
%   - Format: \item \textbf{Category:} Skill1, Skill2, Skill3
%
% SELECTION PRIORITY:
%   1. Match skills to {{JOB_KEYWORDS}} first
%   2. Include skills mentioned in {{JOB_DESCRIPTION}}
%   3. Add relevant complementary skills from asset library
%   4. Order categories by relevance (most important first)
%
% RULES:
%   - Languages section is FIXED - do NOT modify
%   - Keep skill lists concise (3-8 items per category)
%   - Use official tool/language names with correct capitalization
%   - No proficiency levels needed (presence implies proficiency)
%   - If no better title needed, keep default "Technical \& Computational Skills"
% ================================================================

% ================= Skills =================
\cvsection{Skills}

\cvitem{{{SKILLS_SECTION_TITLE}}}{
  \begin{cvsubitems}
    {{TECHNICAL_SKILLS_BULLETS}}
  \end{cvsubitems}
}

\cvitem{Languages}{
  \begin{cvsubitems}
    \item German (Native), English (Professional), French (Working Proficiency)
    % FIXED - Do NOT modify this line
  \end{cvsubitems}
}